\section{The side branches: the Jacobian map, Navier--Stokes and \texorpdfstring{$S^6$}{S6}}\label{sec:side}

Three branches of the record bring material from outside lattice gauge theory into the Hamiltonian setting, each through an explicit typed map. None of them enters the proofs of Sections~\ref{sec:gap} and~\ref{sec:heat}. The attribution file of 21 September says so for the 21 September heat and volume proofs. For the fifth-reference results it follows from their proofs above, which use only the Hamiltonian, its Fourier calculus and the computed coefficients. This section states what each branch proves in this record and how far it was checked. The statements and statuses come from the reader of four workbenches \cite{WorkbenchReader}, where they were refereed, and from the notes cited.

\subsection{The Jacobian map and the low-mode theorem}

The map
\[
F(x,y,w)=\bigl((1+xy)^3w+y^2(1+xy)(4+3xy),\ y+3x(1+xy)^2w+3xy^2(4+3xy),\ 2x-3x^2y-x^3w\bigr)
\]
has $\det DF=-2$ and three preimages of $(-\frac14,0,0)$, namely $(0,0,-\frac14)$ and $(\pm1,\mp\frac32,\frac{13}2)$ (checked symbolically). The workbench's attribution file identifies it with the counterexample to the Jacobian conjecture announced by Levent Alpöge on 20 July 2026, whose announcement, as the file records, credits Akhil for the question and Fable for the work; it notes that the $S^6$ construction is a different object. That Keller maps give incompressible polynomial flows, and that the curve $\gamma(\tau)=(z^{-1},-3z/2,13z^2/2)$, $z=\sqrt{1-8\tau}$, escapes at $\tau=1/8$, is proved in the zeta reader \cite{ZetaReader} (Lemma 6.6, \emph{Keller maps give incompressible polynomial flows}).
The workbench's chain $F\to$ incompressible flow $\to$ material tensor $\to$ lattice local-energy weights $\to$ trial states ends in its Theorem~14.2, audited in \note{34\_}: it is a valid diagonal argument, conditional on fixed-box weak-coupling limits that the file only sketches, and its energy quotients equal, up to a factor $1\pm1/(10j)$, the lowest free two-gluon colour-singlet energy of an open box of side $j/50$ at couplings $g_j<1/j$. So the map carries the Jacobian flow to trial states whose energies track the free two-gluon threshold, and the data of $F$ enter the quotients within the window $1\pm1/(10j)$; which properties of $F$ would move the quotients outside it is Question~\ref{q:jacwindow}. The Jacobian map itself is the subject of the companion paper \cite{JacobianPaper}.

\subsection{Navier--Stokes}

In this record the Navier--Stokes material enters through the quantum line's §§20--22: the velocity field of a solution becomes a reducible $SU(2)$ connection $A_i=\lambda u_iT$, with curvature density $-2\sum_{i<j}\operatorname{tr}F_{ij}^2=\lambda^2|\operatorname{curl}u|^2$, and gauge-invariant curvature states built from it are followed along a weak-coupling diagonal.
\begin{itemize}
\item \emph{The imported theorem} (not audited). The source is OpenAI's manuscript \emph{Finite Time Blowup for Navier--Stokes} \cite{OpenAINS}, Theorem~1.1: for every viscosity, a compactly supported smooth force and zero initial velocity give a smooth solution on $[0,1)$ with bounded energy and $\limsup\|u(t)\|_\infty=\infty$. The record's first \texttt{README.md} says that source statements and targeted derivations were inspected, and that no full independent analytic or Lean verification of the external existence proof is claimed. Every use of it is conditional on it.
\item \emph{The map to Yang--Mills} (audited, conditional; \note{21\_} §2). The identity above and the constants of the curvature masses were checked. The zero-quotient diagonal exists only along couplings $g\to0$, and it rests on rate bounds for the Navier--Stokes solution and a dissipation bound. The companion Navier--Stokes paper derives the rate bounds, conditionally on four named statements of the manuscript \cite[Theorem~3.1]{NSPaper}; the dissipation bound was not checked here. The source calls it ``a rigorous weak-coupling counterexample candidate, not a completed Millennium disproof''.
\end{itemize}
The workbench's separate Navier--Stokes collection (the repository's \texttt{navier-stokes/}, with its own Zenodo record 10.5281/zenodo.22678406) is not part of this record; it is the subject of the companion Navier--Stokes paper \cite{NSPaper}.

\subsection{\texorpdfstring{$S^6$}{S6}}

The $S^6$ construction is the manuscript \cite{AlpogeS6}: a compact complex threefold fibred by tori over the projective line, which the manuscript states is diffeomorphic to $S^6$. That statement was not checked here, and the workbench does not certify a global complex structure on $S^6$; it draws no mass-gap conclusion. Its attribution file, which it designates as the clarification, says that the manuscript was circulated by Levent Alpöge and produced with Claude (older files say \emph{with Fable}).

In this record the construction enters only through its regular-fibre period matrix. The spatial and volume lines use its imaginary-period block to build a magnetic background, magnetic link transports and a gauge-covariant map into the Wilson Hilbert space (the attribution file's \emph{period-dependent magnetic branch}; volume line §1 and §19). On the patch that these lines use, only the combination $D=L_Tq_T+6m_T^2$ of the imaginary periods enters, with $L_T=\operatorname{Im}\tau_T$, $q_T=-\operatorname{Im}\beta_T$ and $m_T=\operatorname{Im}\mu_T$ (\note{28\_}, Lemma~28.2), through the angle $\theta=2\pi a^2/D$. So the Yang--Mills calculations need one positive number from the $S^6$ data, and nothing about $S^6$ itself; the attribution file says the same about its hypotheses.

The $S^6$ project's own results, and the record in which they are published (10.5281/zenodo.22678442), are the subject of the companion $S^6$ paper \cite{SixSpherePaper}, which also determines the monodromy group of the same period system.
